\documentclass[10pt]{article}

\usepackage[a4paper,margin=25mm]{geometry}
\usepackage{booktabs}
\usepackage{microtype}
\usepackage{hyperref}
\usepackage{xcolor}

\hypersetup{
  colorlinks=true,
  linkcolor=blue!55!black,
  citecolor=blue!55!black,
  urlcolor=blue!55!black
}

\title{When 0.99 AUC Is Not an Early-Warning System:\\
A Leakage and Operational-Utility Audit of\\
League of Legends Event Forecasting}
\author{Keerthivasan Kannan}
\date{Working paper --- 15 September 2026}

\begin{document}
\maketitle

\begin{abstract}
Short-horizon event prediction in multiplayer online battle arena games is
usually reported as row-level classification, although a deployed warning
system acts on distinct events and can burden users with repeated false alerts.
We audit a released League of Legends early-warning artifact containing
484,255 timestamped rows from 2,998 matches. The original pipeline creates
overlapping windows, makes temporal-jitter copies, and randomly splits the
result. Reconstructing this procedure shows that 99.10\% of test examples have
the same base-window family in another partition, 100\% share a source match,
and 100\% overlap another partition in time. Sixteen final match-summary fields
also appear at every timestamp, eleven intended live-state fields are constant,
and the release lacks a patch identifier. We replace the evaluation with
complete-match splits, causal feature exclusions, validation-only thresholds,
upward-crossing alerts, cooldowns, and one-to-one event matching. A causal
gradient-boosted baseline reaches row-level ROC--AUCs of 0.990, 0.959, and 0.743
for Baron, Dragon, and a teamfight proxy, respectively, but event-level alert
F1 scores are only 0.376, 0.415, and 0.149. Teamfight warnings generate 3.37
false alerts per game. These results do not establish a new predictive model;
they establish that discrimination, temporal independence, calibration, and
operational utility must be evaluated separately. We release executable audit
and benchmark code and preregister a future-patch research protocol.
\end{abstract}

\section{Introduction}

Forecasting an objective fight or a multi-player engagement before it occurs
could support replay review, coaching, and broadcast analysis. The same task is
also unusually vulnerable to misleading evaluation. Objective availability is
strongly tied to match time, observations from one match are highly dependent,
future labels occupy several adjacent rows, and overlapping sequence windows
share most of their input. A conventional random row split can therefore test
recognition of a familiar match or window rather than generalization.

Prior work has predicted several important events in Honor of Kings and studied
prediction attribution \cite{yang2022events}; recent work also considers MOBA
event prediction \cite{vardakis2026moba}. League-specific graph measures have
been used for aggregate outcome and collective-intelligence analysis
\cite{caldeira2025collective}, while patch-agnostic representations have been
investigated in Dota 2 \cite{chitayat2023meta}. These precedents mean that an
attention layer, graph representation, or event target is not by itself a novel
contribution. Evaluation must demonstrate information beyond schedule priors
and usefulness beyond a row-level score.

This paper makes three bounded contributions. First, it provides an executable
forensic audit of a public educational artifact. Second, it contrasts
row-level discrimination with event-level warning utility under a stricter
protocol. Third, it preregisters the evidence required for a future temporal
graph and discrete-hazard study. We make no claim that the audited data support
future-patch or player-facing deployment.

\section{Artifact and audit method}

The audited CSV is identified by a full SHA--256 digest in the public artifact
manifest. It contains 86 columns and a regular ten-second time index. We traced every
column to the released generation notebook and every split to the released
model notebook, then encoded structural checks as tests and a command-line
audit.

\subsection{Information availability}

Sixteen team aggregate fields are computed from end-of-match participant
summaries and merged back onto every timestamp. They include final kills,
deaths, assists, damage, healing, and ward counts. These values are unavailable
at prediction time. Eleven further fields are constant: alive counts, low-health
counts, and objective counts. The generator asks timeline participant frames
for a health field they do not supply; the default zero makes alive counts zero
and low-health counts five. Objective attribution reads a generic team field
rather than the objective event's killer-team field.

The regular ten-second table overstates source resolution. Gold and experience
snapshots change in only 16.7\% of rows because lower-frequency frames are
forward-filled. The teamfight proxy adds each kill that begins a ten-second
interval containing at least three kills; a single combat episode can therefore
produce several event timestamps. Finally, no creation time or game version is
retained, preventing a true chronological or future-patch test.

\subsection{Split reconstruction}

The original model code builds length-40 windows every five rows, producing
74,728 base windows. It creates two jittered copies plus the original, yielding
224,184 examples, and then applies random 90/10 and 80/20 splits. We reproduced
the exact index operations with random seed 42 and retained source match,
base-window family, and nominal interval metadata. Table~\ref{tab:contamination}
measures whether each example has related material in another partition.

\begin{table}[ht]
\centering
\caption{Cross-partition contamination in the released split procedure.}
\label{tab:contamination}
\begin{tabular}{lrrrr}
\toprule
Partition & Examples & Same family & Same match & Time overlap \\
\midrule
Train      & 161,412 & 48.47\% & 100.00\% & 100.00\%$^{*}$ \\
Validation &  40,353 & 96.83\% & 100.00\% & 100.00\% \\
Test       &  22,419 & 99.10\% & 100.00\% & 100.00\% \\
\bottomrule
\multicolumn{5}{l}{\footnotesize $^{*}$99.998\% before rounding.}
\end{tabular}
\end{table}

Global standardization is also fitted before splitting. Consequently, the
reported legacy test metrics are not estimates of performance on unseen
matches.

\section{Leakage-safe baseline protocol}

Because creation timestamps are absent, we sort by the numeric part of match
ID as an explicitly imperfect chronology proxy. Complete matches are allocated
70/15/15, giving 2,098 training, 450 validation, and 450 test matches. No match
crosses partitions. Transformations are fitted after splitting.

We evaluate three deterministic histogram gradient-boosted baselines: match
time only; time plus prior event/history features; and 49 numeric causal
features after removing labels, identifiers, constant generator failures, and
post-match summaries. The target is an event in the next 30 seconds. This
legacy target reconstruction is retained only for comparison and is not the
planned research-v2 label source.

An operating threshold is chosen by maximum event-level F1 on validation
matches only. An alert occurs on an upward threshold crossing and triggers a
60-second cooldown. Greedy chronological matching allows each alert and event
to be used at most once when the alert precedes the event by at most 30 seconds.
We report average precision because events are rare, Brier score as an initial
probability diagnostic, and operational precision, recall, F1, false alerts per
game, and lead time. Class balancing means these raw scores are not claimed to
be calibrated; calibration is a separate preregistered stage.

\section{Results}

Table~\ref{tab:results} reports the untouched test partition. High Baron
ROC--AUC is already achieved with match time alone (0.960), demonstrating a
large schedule prior. History and causal state improve average precision and
alert utility, but no event type satisfies the preregistered product gate of
precision and recall at least 0.50. The causal Dragon model reaches precision
0.515 but recall 0.348. The teamfight proxy is particularly weak and produces
more than three false alerts per game.

\begin{table}[ht]
\centering
\caption{Test discrimination and operational warning results.}
\label{tab:results}
\small
\begin{tabular}{llrrrrrr}
\toprule
Features & Event & ROC--AUC & AP & Precision & Recall & Alert F1 & False/game \\
\midrule
Time    & Baron     & .960 & .143 & .213 & .225 & .219 & .49 \\
History & Baron     & .989 & .487 & .316 & .419 & .361 & .54 \\
Causal  & Baron     & .990 & .494 & .344 & .416 & .376 & .47 \\
Time    & Dragon    & .672 & .101 & .097 & .125 & .109 & 3.26 \\
History & Dragon    & .954 & .526 & .440 & .369 & .401 & 1.31 \\
Causal  & Dragon    & .959 & .562 & .515 & .348 & .415 & .91 \\
Time    & Teamfight & .674 & .171 & .086 & .089 & .087 & 2.39 \\
History & Teamfight & .740 & .266 & .123 & .165 & .141 & 2.97 \\
Causal  & Teamfight & .743 & .271 & .123 & .188 & .149 & 3.37 \\
\bottomrule
\end{tabular}
\end{table}

Median lead time is 20 seconds for all causal models. A deliberately invalid
diagnostic trained only on the 16 post-match summary fields produces no useful
test alerts. This is a negative result: unavailable summaries are still
leakage, but final totals alone do not localize the next event.

\section{Implications for event-forecasting research}

The gap between 0.990 row ROC--AUC and 0.376 alert F1 is not a contradiction.
ROC--AUC evaluates ranked rows across all thresholds; the warning system emits
distinct interventions in a dependent time stream. Many easy negative rows can
support a high ROC--AUC while repeated or mistimed crossings remain unhelpful.
Schedule priors make this especially severe for Baron.

A credible successor should preserve exact event timestamps and native frame
cadence, split complete matches before preprocessing, reserve future patches,
and parameterize several horizons coherently. Temporal Graph Networks provide a
general representation for timed interaction graphs \cite{rossi2020tgn}, but a
graph model should be accepted only if it beats clock, history, tabular, and
sequence baselines under paired match-level inference. Patch rules should be
inputs rather than hard-coded constants.

Our registered successor models players and objectives as dynamic nodes and
uses discrete hazards so cumulative risk is monotone across 10, 20, 30, and 60
seconds. It tests future-patch conditioning and OOD abstention. These are
hypotheses, not results.

\section{Limitations}

The released table does not permit validation against raw event JSON, a true
time split, patch analysis, or a strict repeat-player holdout. Match ID ordering
is only a proxy. The teamfight label is a behavioral proxy rather than an
official semantic event. Our alert matcher and cooldown encode one plausible
user experience; sensitivity analysis is needed. We have not yet produced
match-bootstrap confidence intervals. The recent Vardakis et al. article must
be reviewed in full before final novelty claims or submission.

\section{Conclusion}

The audited model results should not be used as evidence for deployment. Their
split allows almost every test example to share an augmented family and every
test example to share a match and overlapping time with another partition.
After enforcing match independence and evaluating distinct warnings, strong
row discrimination does not translate into adequate alert utility. This
negative finding provides a useful foundation: future work should earn its
complexity through exact data provenance, future-patch generalization,
calibration, alert burden, and falsifiable comparison with simple priors.

\section*{Reproducibility statement}

The repository records the dataset hash, audit findings, split policy and hash,
runtime versions, model configuration, validation thresholds, and
machine-readable test results. Raw and derived datasets are distributed
separately subject to their licences. The source code is MIT licensed. AI use
is disclosed in the repository; a human author remains responsible for every
claim and reference.

\bibliographystyle{plain}
\bibliography{references}

\end{document}
